\roadmapSection{Book}{DATA IS DNA --- Data and Computation Foundations}{Write During Section 0, Then Ongoing}

{\small\itshape\color{arligray} Data Is DNA asks the question every other book skips: what IS data? How does information become computation? Write early --- this book integrates everything you learn.}

\textbf{Repository:} \texttt{github.com/papyrxis/data-is-dna}

\roadmapHead{Writing Discipline}
\checkitem{Start writing Chapter 1 in Month 1 --- it maps directly to your electronics fundamentals work}{}
\checkitem{Every hardware and software concept you learn has a Data Is DNA chapter. Write the connection same day}{}
\checkitem{This book is deliberately broad --- each chapter is a 2000-word essay, not a textbook chapter}{}
\checkitem{One chapter draft per month minimum. Revise every two months}{}

\vspace{0.3cm}

% ── Chapter 1 ────────────────────────────────────────────────────────────────
\roadmapSubSection{D.1}{Chapter 1: The Nature of Information and Reality}{Write in Month 1}

{\small\itshape\color{arligray} The philosophical foundation. Connects to Shannon information theory and your electronics work.}

\checkitem{What Is Information? --- Shannon entropy, bits as uncertainty reduction}{}
\checkitem{Data, Meaning, and Representation --- the gap between physical signals and meaning}{}
\checkitem{Patterns, Symbols, and Abstraction --- how humans impose structure on physical phenomena}{}
\checkitem{Entropy and Uncertainty --- quantify information: H = −Σ p log₂(p)}{}
\checkitem{Analog vs Digital Reality --- why we discretize, what we lose, what we gain}{}
\checkitem{Connect to your work: every repair you do is working with information encoded in voltage levels}{}

% ── Chapter 2 ────────────────────────────────────────────────────────────────
\roadmapSubSection{D.2}{Chapter 2: Mathematical Foundations of Computation}{Write in Month 2}

\checkitem{Binary as Universal Representation --- why base-2, why not base-10}{}
\checkitem{Boolean Algebra --- gates as math. Connect to Section 14.5 digital logic work}{}
\checkitem{Sets, Relations, and Functions --- mathematical formalism of data structures}{}
\checkitem{Discrete Mathematics overview --- connect to Mathesis Part V}{}
\checkitem{Matrices, Vectors --- linear algebra as data transformation}{}

% ── Chapter 3 ────────────────────────────────────────────────────────────────
\roadmapSubSection{D.3}{Chapter 3: From Physics to Electronics}{Write in Month 3}

{\small\itshape\color{arligray} This chapter is your electronics section written as a book chapter. Write it while doing Section 14.}

\checkitem{Voltage, Current, Resistance, Power --- map Ohm's Law to information theory}{}
\checkitem{Transistors and Switching --- how a physical device becomes a logical 1 or 0}{}
\checkitem{Digital Signals and Logic Levels --- voltage thresholds, noise margins, the abstraction layer}{}
\checkitem{Logic Gates and Circuits --- from MOSFET to NAND to full adder}{}
\checkitem{Clock Signals and Synchronization --- time as the organizer of digital information}{}
\checkitem{Physical Limits of Computation --- quantum tunneling, heat, why Moore's Law slows}{}

\newpage
% ── Chapter 4 ────────────────────────────────────────────────────────────────
\roadmapSubSection{D.4}{Chapter 4: Computer Architecture}{Write in Month 5--6}

\checkitem{Von Neumann Model --- stored-program computer, the same architecture in every phone you repair}{}
\checkitem{Processors and Execution Units --- connect to ARLIZ Vol II CPU pipeline}{}
\checkitem{Memory Hierarchies --- registers → L1 → L2 → RAM → storage. Each level a different data model}{}
\checkitem{Input and Output Systems --- buses, interrupts, DMA. Connect to your embedded work}{}
\checkitem{Parallelism and Modern Architectures --- multi-core, SIMD, GPUs}{}

% ── Chapters 5–8 ─────────────────────────────────────────────────────────────
\roadmapSubSection{D.5}{Chapters 5–8: Memory, Data Structures, Algorithms, Type Systems}{Write Months 6--12}

\checkitem{Ch.5 Memory and Representation: bits, bytes, IEEE 754, ASCII, endianness --- same as ARLIZ Vol I Ch.45--64}{}
\checkitem{Ch.5 Pointers, alignment, compression --- write after completing Section 4.1 C programming}{}
\checkitem{Ch.6 Data Structures --- why they exist, trade-offs. Write in parallel with ARLIZ Vol III}{}
\checkitem{Ch.7 Algorithms --- complexity, searching, sorting. Write alongside Art of Algorithmic Analysis}{}
\checkitem{Ch.8 Type Systems --- static vs dynamic, generic programming, type safety. Write after C section}{}

% ── Chapters 9–10 ────────────────────────────────────────────────────────────
\roadmapSubSection{D.6}{Chapters 9–10: Systems and Future of Computation}{Write Months 13--18}

\checkitem{Ch.9 OS, Networks, Databases, Distributed Systems --- write after Estonia prep research}{}
\checkitem{Ch.9 Security and Cryptography --- write alongside Section 6 Reverse Engineering}{}
\checkitem{Ch.10 AI and Neural Networks as data transformation systems}{}
\checkitem{Ch.10 Quantum Computing, Biological Computation --- the frontier}{}
\checkitem{Ch.10 The Philosophy of Intelligence --- close the book where it opened: what IS information?}{}
